% problem-000095 1 配位化学与氧化态
\section*{1. 配位化学与氧化态}
\textit{下列为分问。}

\subsection*{1-1}
电⼦⼯业中硅器件的腐蚀可以使⽤重铬酸盐溶液（反应1）或者氢氟酸和浓硝酸的混合液（反应2），写出反应1的离⼦反应⽅程式和反应2的化学反应⽅程式。

\subsection*{1-2}
硅酸盐的结构和性质研究是⽆机化学内⼀个重要的⽅向。

\subsection*{1-2}
-1 硅酸盐基本结构单元 $\mathrm { [ S i O _ { 4 } ] ^ { 4 } }$ −中的Si—O平均键⻓为162 pm，从Si、O原⼦间的成键特征⻆度解释为何该平均键⻓⼩于Si和O原⼦的共价单键半径之和(190pm)。

\subsection*{1-2}
-2 上述硅酸盐结构中的Si(IV)被Al(III)部分同构取代后会产⽣Brönsted（布朗斯特）酸位，有如下图所⽰的模型A和B被提出来解释该酸性的产⽣。哪种模型结构中的羟基更易脱去质⼦？原因是什么？

（图：）
模型A

（图：）
模型B

\subsection*{1-3}
硅烷是⼀类重要的有机硅化合物，以π型官能化硅烷和杂原⼦型官能化硅烷为配体合成得到的配位化合物在催化和光电领域有可观的应⽤。

\subsection*{1-3}
-1 研究表明，以硅原⼦为桥的⼆茂铁配合物 $\mathrm { [ F e ( n – C _ { 5 } H _ { 4 } ) _ { 2 } S i R _ { 2 } ] }$ （R为烷基）可发⽣开环聚合⽽⽣成配位聚合物。画出R为甲基时该配合物和相应聚合产物的结构 $\updownarrow \updownarrow _ { \updownarrow }$

\subsection*{1-3}
-2 配合物 $\mathrm { { [ F e ( n – C _ { 5 } H _ { 4 } ) _ { 2 } S i R _ { 2 } } }$ ]的R基团为⼄基时⽣成的配位聚合物记为A，R基团为正丁基时⽣成的配位聚合物记为B；A和B的聚合度相同。A和B中哪⼀种具有较低的玻璃化转变温度？简述理由。

\subsection*{1-3}
-3 如下图所⽰的杂原⼦官能化硅烯在正已烷溶液中和羰基配合物 $\mathrm { N i ( C O ) _ { 4 } }$ _{4}，发⽣配体取代反应释放出等物质的量的 $\mathrm { C O } ,$ ，⽣成化合物 $\mathrm { c } ;$ ；C继续和该硅烯发⽣配体取代反应释放出等物质的量的 $\mathrm { C O } _ { \mathrm { { z } } }$ ，⽣成化合物D；上述反应前后配合物的空间构型不变，画出化合物C和D的结构式。

（图：）
